% titel: Blatt 9, A4: Osmotischer Druck von Glukose
% Teil II: der durchgerechnete Weg.

\subsection*{Gegeben und gesucht}
6.00 g Glukose (C$_6$H$_{12}$O$_6$, $M=180.16$) in 100 mL Wasser bei
25.0\,°C. Gesucht: $\Pi$.

\subsection*{Rechnung}
\[
  n = \frac{6.00}{180.16} = 0.03331\ mol,\qquad
  c = \frac{0.03331}{0.100} = 0.3331\ mol\,L^{-1}.
\]
\[
  T = 25.0 + 273.15 = 298.15\ K,\qquad
  \Pi = c\,R\,T.
\]
Mit $R = 0.08314\ L\,bar\,mol^{-1}K^{-1}$:
\[
  \Pi = 0.3331\cdot0.08314\cdot298.15 \approx 8.26\ bar .
\]

\begin{ergebnis}
$\Pi \approx 8.3\ bar$ (Glukose, 6.00 g in 100 mL, 25\,°C).
\end{ergebnis}

\begin{kontrolle}
Andere Einheit: $\Pi=cRT
=333.1\ mol\,m^{-3}\cdot8.314\cdot298.15 \approx 8.26\cdot10^5\ Pa
\approx 8.3\ bar$. Konsistent.
\end{kontrolle}

\begin{merksatz}
$\Pi = cRT$ --- erst $c=n/V$ und $T$ in Kelvin; dann passende
Gaskonstante wählen.
\end{merksatz}
